\chapter{Workspace characterisation and its consequences}
\label{ch:workspace}

This chapter is separated from \cref{ch:results} because it is not a single
measurement but a chain of them, and because its consequences determined the
shape of the task set. It reports a negative finding together with a
quantified and buildable mechanical fix; the finding alone would be a
complaint rather than a contribution.

\section{Continuous reachability from the home pose}

Reachability was measured as a swept extent rather than as pointwise
feasibility, because teleoperation requires the arm to travel to a pose from
where it already is, without a branch change. Twenty-six directions were swept
in \SI{1}{\centi\metre} steps from the home configuration, each step seeded
from the previous solution exactly as the follower does.

Repeating the sweep ten times gave identical quartiles: the inter-quartile
range was \SI{0.000}{\metre} in every direction for both arms. The measurement
repeats. What did \emph{not} repeat in earlier runs was the summary statistic:
isotropy taken as the minimum over directions of a single sweep is dominated
by rare early terminations and varied by a factor of two between runs. On
medians it is stable, at \num{0.16} for the left arm and \num{0.10} for the
right (\cref{fig:reach}).

\begin{figure}[htbp]
  \centering
  \includegraphics[width=0.95\textwidth]{reach_directions.pdf}
  \caption[Continuous reachability from home]{Continuous reachability from the
  home pose in the six principal directions, median of ten sweeps, with the
  binding constraint printed inside each bar. The envelope is markedly
  anisotropic: the median reach over all twenty-six directions is
  \SIrange{0.34}{0.40}{\metre} while the worst direction is
  \SI{0.14}{\metre} on the left arm and \SI{0.09}{\metre} on the right.
  \emph{Max range} marks a direction in which the sweep exhausted its own
  travel before the arm did, so those two figures are lower bounds.
  Regenerated from \repo{recordings/baselines/workspace\_n10\_20260806.json}.}
  \label{fig:reach}
\end{figure}

The early terminations were traced to the measuring tool rather than to the
robot. A sweep is a chain of up to ninety solver calls and any single failure
ends that direction; the follower does not suffer this because it re-seeds the
redundant joint and retries. Giving the sweep the same retry reduced early
terminations from nine to zero on the left arm and three to zero on the right,
while the medians barely moved. That the medians were already right is what
establishes the outliers as chain compounding rather than geometry.


\begin{figure}[htbp]
  \centering
  \includegraphics[width=0.72\linewidth]{rviz/home_pose_labelled.png}
  \caption[The mounted configuration at the home pose]{The simulated platform
  at the home joint angles, captured from the running system rather than
  drawn. Both arms are parked over the shoulders and both hands finish on the
  same side of the body, \SI{1.46}{\metre} apart. That asymmetry is not a
  rendering artefact and not a mount error: it follows from the home joint
  angles, which are ground truth because the simulation-to-hardware bridge
  replays them onto the real arm. It is the disjoint-workspace result of this
  chapter made visible.}
  \label{fig:rviz-home-workspace}
\end{figure}

\section{The two arms share no reachable workspace}

Over a grid of sixty-three cells spanning the frontal working volume, with the
arms verified at the home pose, \textbf{no cell was reachable by both arms}.
The result survives the two obvious objections, as \cref{tab:overlap} shows:
freeing the wrist yaw raises each arm's individual reach from sixteen and eighteen
cells to twenty-one and still yields no shared cell, and repeating the test
with more redundancy retries changes nothing.

\begin{table}[htbp]
  \centering
  \caption[Both-arm overlap under three test conditions]{Both-arm overlap over
  the 63-cell frontal grid, arms verified at home. The null result survives
  both wrist freedom and additional inverse-kinematics retries.}
  \label{tab:overlap}
  \begin{tabular}{lccc}
    \toprule
    Test condition & Both & Left & Right \\
    \midrule
    Fixed orientation, 3 retries & \textbf{0/63} & 16 & 18 \\
    Fixed orientation, 6 retries, 3 repeats & \textbf{0/63} & 17 & 18 \\
    Yaw free, 6 retries & \textbf{0/63} & 21 & 21 \\
    \bottomrule
  \end{tabular}
\end{table}

A fine sweep in \SI{0.05}{\metre} steps across the plane on which the arms
would have to meet locates the boundary precisely: the right arm reaches
$x \le \SI{-0.10}{\metre}$, the left reaches $x \ge \SI{0.15}{\metre}$, and
there is a \SI{0.25}{\metre} dead band between them in which neither arm can
place its end effector (\cref{fig:disjoint}).

\begin{figure}[htbp]
  \centering
  \includegraphics[width=0.9\textwidth]{workspace_disjoint.pdf}
  \caption[The measured dead band]{The two reachable sets do not meet. Measured
  at $y=\SI{0.35}{\metre}$, $z=\SI{1.10}{\metre}$ with yaw free and three
  repeats per point.}
  \label{fig:disjoint}
\end{figure}

\subsection{What this does and does not establish}

It establishes that, within $x \in [-0.6, 0.6]$, $y \in [0, 0.4]$,
$z \in [0.90, 1.30]$~\si{\metre}, with the as-built mount and the legacy home
joint angles, the two reachable sets do not intersect.

It does \emph{not} establish behaviour outside that volume, nor does it
generalise to other mounts. The correct claim is that the arms share no
workspace \emph{as mounted and as parked}, and the remedy is mechanical.

\section{The fix is buildable}

A parametric sweep over candidate mount placements shows the null result is
not intrinsic. Moving the mounts inboard, forward and down, with additional
forward tilt, raises the number of shared cells from three to
\textbf{twenty-two} on the purely kinematic measure, within a buildable
\SI{0.20}{\metre} mount separation (\cref{fig:mount}).

\begin{figure}[htbp]
  \centering
  \includegraphics[width=0.85\textwidth]{mount_sweep.pdf}
  \caption[Mount parameter sweep]{Single-parameter sweeps from the as-built
  mount. Forward tilt and inboard translation both help substantially; splaying
  the mounts outward \emph{hurts}. The sweep is purely kinematic, moving the
  arms' bases without moving the wearer, so these are upper bounds and any
  candidate requires a clearance check before adoption.}
  \label{fig:mount}
\end{figure}

The sweep was recorded and deliberately not applied at the time, because
changing the mount invalidates the home pose, the workspace anchor and every
clearance figure reported in \cref{ch:results}. It was applied later, as one
deliberate pass with the whole task set re-measured afterwards, and
\cref{sec:mount-applied} reports what it bought.

\section{One boundary, two independent constraints}

Two further measurements, taken for different reasons, land on the same
lateral boundary and together make the case that the inboard band is
unusable.

The first is wearer clearance. Every scenario in the chest-height band passes
inverse kinematics with collision checking enabled, and the planner reports
the poses valid, yet the arms pass within \SIrange{48}{67}{\milli\metre} of
the wearer's torso, inside the margin the real-robot configuration enforces,
which was \SI{120}{\milli\metre} at the time of that measurement and is
\SI{150}{\milli\metre} now. The reconciliation is that the two tests ask different
questions: the planner's check is binary contact, whereas the floor is a
margin, and a pose can satisfy one while being refused by the other. Sweeping
the lateral half-span shows clearance rising monotonically and crossing the
floor at \SI{0.25}{\metre}.

The second is grasp feasibility. A top-down grasp requires the wrist to adopt
an orientation with its approach axis vertical. Sweeping the same lateral
coordinate gives a sharp cliff: zero of nine grasp candidates are solvable at
$|x| = \SI{0.25}{\metre}$ and nine of nine at $|x| = \SI{0.30}{\metre}$. The
failure persists with collision checking disabled, so it is kinematic
infeasibility of the wrist orientation rather than obstruction.

\begin{figure}[htbp]
  \centering
  \includegraphics[width=0.85\textwidth]{lateral_constraint.pdf}
  \caption[Two constraints, one boundary]{Wearer clearance and top-down grasp
  feasibility against lateral half-span. The two constraints were measured
  independently and for different purposes, and both place the usable boundary
  at $|x| \approx \SIrange{0.25}{0.30}{\metre}$.}
  \label{fig:lateral-workspace}
\end{figure}

\section{What the pinned wrist costs, at the point where the work happens}
\label{sec:work-point}

\Cref{sec:sim-orientation} measures reach under each orientation policy by
walking outward from each arm's \emph{home} end-effector pose, and finds that
freeing the wrist buys \SI{87}{\milli\metre} of mean reach. Walking from the
home pose is the natural thing to do and it understates the cost badly,
because home is in free space and the task is not.

Repeating the walk from the point the task actually works at ---
$(\num{0.4}, \num{0.175}, \SI{1.12}{\metre})$, on the working surface in front
of the wearer --- and solving through the follower's own orientation policy
rather than against a continuous bound, gives a different picture entirely.

\begin{table}[htbp]
  \centering
  \small
  \caption[Reach from home and from the work point]{Mean reach under each
  orientation policy, from two start points, both with the
  \SI{150}{\milli\metre} wearer floor enforced throughout. The two columns are
  the same measurement of the same arms and they disagree by a factor of
  seven, because a reach figure is meaningless without the point it was
  measured from.}
  \label{tab:work-point}
  \begin{tabular}{@{}lcc@{}}
    \toprule
    Orientation policy & from \textsc{home} & from the \textsc{work point} \\
    & \si{\metre} & \si{\metre} \\
    \midrule
    pinned to the anchor (delivered) & \num{0.460} & \textbf{\num{0.065}} \\
    roll free                        & \num{0.463} & \num{0.065} \\
    within a \SI{15}{\degree} cone   & \num{0.492} & \num{0.304} \\
    within a \SI{45}{\degree} cone   & \num{0.530} & \num{0.325} \\
    position only                    & \num{0.547} & \num{0.363} \\
    \bottomrule
  \end{tabular}
\end{table}

\textbf{At the point where the work happens, the delivered policy leaves the
arm \SI{65}{\milli\metre} of mean reach against \SI{363}{\milli\metre} with
the wrist free --- a factor of \num{5.6}.} That is the honest size of the
constraint this system operates under, and it is why the arms appear barely to
move while every reachability figure taken at home looks comfortable.

Three things follow.

\emph{The mechanism is redundancy, not the wearer.} No walk from the work
point stops because of the clearance floor; they stop because there is no
solution. A fixed six-degree-of-freedom pose on a seven-degree-of-freedom arm
spends the entire redundancy, and the redundancy is the joint that moves the
elbow out of the way while the hand stays put. Near the body that joint is the
only thing that makes a pose reachable at all, and pinning the wrist removes
it.

\emph{Freeing the roll alone buys nothing, in either column.} It is worth
\SI{3}{\milli\metre} from home and \SI{0}{\milli\metre} from the work point.
The redundancy does not live in the roll, so the cheapest-sounding relaxation
is the one that is worth least --- which is a useful thing to know before
spending a session on it.

\emph{Almost all of the recovery comes from the first \SI{15}{\degree}.} A
\SI{15}{\degree} cone recovers \SI{239}{\milli\metre} of the
\SI{298}{\milli\metre} that separates pinned from free, or
\SI{80}{\percent} of it. The follower already supplies a \SI{15}{\degree}
cone when it plans, which is why the task set solves at all
(\cref{sec:sim-tasks}) --- and it means the delivered system is not in fact
operating in the \SI{0.065}{\metre} column. The pinned figure is what the
\emph{commanded} orientation would cost if the follower did not quietly widen
it, and stating that separation is the point of the table.

\section{The centre of the working space, asked four ways}
\label{sec:centre}

The question the platform is most often asked is whether the arms can work in
front of the wearer's own chest, where the wearer's own hands would be. It was
asked four times, by four instruments built for different purposes, because
the first answer was unwelcome enough to distrust.

\begin{table}[htbp]
  \centering
  \caption[The centre, four measurements]{Four independent attempts to place
  usable work within \SI{100}{\milli\metre} of the body centre line. All are
  at ten repeats over the full densified path with the wearer measured
  geometrically against the \SI{150}{\milli\metre} floor.}
  \label{tab:centre}
  \begin{tabular}{@{}p{0.34\linewidth}p{0.58\linewidth}@{}}
    \toprule
    Question & Answer \\
    \midrule
    Innermost workable column at the working height &
      $|x| = \SI{0.325}{\metre}$ on the left arm and \SI{0.450}{\metre} on the
      right \\
    The same, with the wearer's arms deleted from the model entirely &
      \SI{0.300}{\metre} and \SI{0.375}{\metre}. Every column from
      \SI{0.125}{\metre} to \SI{0.300}{\metre} then binds on the
      \textbf{torso}, and at $|x| = \SI{0.10}{\metre}$ the clearance is
      \SI{-0.007}{\metre} and \SI{-0.013}{\metre}: the arm is inside the chest \\
    Objects \emph{resting} on a surface, sweeping surface height and distance
      together, \num{3360} configurations &
      \textbf{no configuration} places a graspable object within
      \SI{100}{\milli\metre} of the centre line \\
    The same, asked for a top-down approach instead of the side approach &
      \textbf{0 of 840} cells, at any surface height from \SI{0.70}{\metre} to
      \SI{1.10}{\metre}, on either arm \\
    \bottomrule
  \end{tabular}
\end{table}

Two conclusions follow, and the second is the operationally useful one.

The constraint is the \textbf{torso}, not the wearer's limbs. Making the
wearer's arm posture a variable and re-running the sweep across five postures
moves the boundary by at most \SI{75}{\milli\metre}, and deleting the arms
altogether leaves the centre shut. An instruction to the wearer is therefore
worth something and is not a route to the centre. Two of the five postures are
actively worse than the default: with a back-mounted rig, arms clasped behind
the back and arms held out to the sides both place the wearer's own limbs
inside the clearance floor at the home pose, before the robot has been asked
to do anything.

\textbf{Neither arm crosses the centre line at all.} Measured at ten repeats
with the wearer and the working surface present, every cross-side placement
slot and every cross-side object returns zero solutions of ten, on both arms.
Not marginal: nothing solves. This is what forces the randomised stage of the
pick-and-place task to draw the \emph{side} and let the colour follow, and the
task builder refuses an object whose declared colour names the other side's
target rather than emitting a path whose every waypoint fails for a reason
nobody would connect to colour.

\section{What the applied mount change bought}
\label{sec:mount-applied}

The mount modification was subsequently applied, as one pass with the task set
re-measured afterwards: the mounts moved \SI{150}{\milli\metre} outboard with
\SI{15}{\degree} of outboard yaw, the working surface was raised, and its near
edge was brought forward.

\begin{table}[htbp]
  \centering
  \caption[Before and after the mount change]{What the applied mount change
  altered. Home clearance improves because the mount is no longer the nearest
  part of the robot to the wearer, which also raises the ceiling on every
  clearance figure the platform can report.}
  \label{tab:mountapplied}
  \begin{tabular}{@{}lll@{}}
    \toprule
    Quantity & Before & After \\
    \midrule
    Clearance at the home pose & \SI{0.1610}{\metre} & \SI{0.2202}{\metre} \\
    Nearest robot part to the wearer at home & the mount, which no joint moves
      & the moving chain \\
    Task objects resting on the surface & not achievable & achieved \\
    Pick-and-place placement targets & \SI{0.530}{\metre} and
      \SI{0.350}{\metre} off centre & \SI{0.290}{\metre} either side \\
    \bottomrule
  \end{tabular}
\end{table}

The cost is that every quantity derived from the previous mount had to be
re-derived, which is the reason the change was deferred rather than taken
early: the approach anchor, the pad offset, the home end-effector position and
every task coordinate computed through them. That is a session's work and it
is the correct way to spend it, but it cannot be done twice.

\section{Clearance along the task paths}

Reachability is not the binding property for a machine mounted on a person.
\Cref{fig:clearance} gives the closest approach to the wearer over the whole
of each task's path, which is the quantity the safety case turns on.

\begin{figure}[htbp]
  \centering
  \includegraphics[width=0.95\textwidth]{clearance_tasks.pdf}
  \caption[Clearance by task]{Closest approach to the wearer over each task's
  full commanded path, ten repeats, with the body measured geometrically and
  the nearest body part printed above each bar. Only the pick-and-place task
  clears the floor on both arms. The three bars at \SI{161}{\milli\metre} are
  the mount before the change of \cref{sec:mount-applied}, which no joint
  moves and which therefore sets the ceiling. Regenerated from
  \repo{recordings/baselines/home\_change\_applied.json}.}
  \label{fig:clearance}
\end{figure}

The breaches are reported rather than corrected. Each is a layout change with
its own re-verification, and moving a task's coordinates to satisfy the floor
without walking the whole path again would replace a known breach with an
unknown one.

\section{Consequences for the task set}

Three consequences follow directly, and each is visible in the task
definitions of \cref{ch:results}.

\textbf{Inter-arm handover is impossible.} It requires both arms at one point,
and no such point exists. This removes an entire task family, and no control,
planning or autonomy change can restore it.

\textbf{A bimanual task must either span the gap or exploit simultaneity.} A
rigid object held at two points can bridge the dead band; two targets in the
two disjoint sets can require simultaneous satisfaction. Both routes are used.

\textbf{The coupled objects had to be re-specified.} The original
\SI{310}{\milli\metre} tray placed each grip at $|x| = \SI{0.155}{\metre}$,
inside the dead band, and ten of twenty-five waypoints failed once the arms
were verified at home. Widening to \SI{500}{\milli\metre} places each grip at
$|x| = \SI{0.25}{\metre}$, which satisfies the clearance boundary and the
reachability boundary simultaneously. The compliant carrier had to lengthen
with it: retention of a ball of radius $r$ requires sag of at least $2r$, so

\begin{equation}
  \mathrm{sag}(s) = \sqrt{\left(\tfrac{L}{2}\right)^{2}
                          - \left(\tfrac{s}{2}\right)^{2}},
  \qquad
  s_{\max} = 2\sqrt{\left(\tfrac{L}{2}\right)^{2} - (2r)^{2}},
  \label{eq:sling}
\end{equation}

where $s$ is the gripper separation, $L$ the carrier length and $r$ the ball
radius. For $L = \SI{540}{\milli\metre}$ and $r = \SI{20}{\milli\metre}$ this
gives $s_{\max} = \SI{534}{\milli\metre}$, leaving a \SI{34}{\milli\metre}
margin at the nominal \SI{500}{\milli\metre} separation
(\cref{fig:sling}). The length is chosen so that the failure threshold sits
\emph{inside} the reachable band: a longer carrier could never fail and would
measure nothing.

\begin{figure}[htbp]
  \centering
  \includegraphics[width=0.8\textwidth]{sling_geometry.pdf}
  \caption[Compliant carrier geometry]{Sag against gripper separation for the
  re-specified compliant carrier. The nonlinearity is the point: sag falls
  slowly, then the ball is lost abruptly at \SI{534}{\milli\metre}.}
  \label{fig:sling}
\end{figure}
